\documentclass{article}
\usepackage[utf8]{inputenc}
\usepackage[margin=1in]{geometry}
\usepackage{amsmath}
\usepackage{amssymb}
\usepackage{braket} % For Dirac notation < | >
\usepackage{fancyhdr}
\usepackage{enumitem}
\usepackage{pgfplots}
\usepackage{tikz}

% Setup for headers and footers
\pagestyle{fancy}
\fancyhf{} % Clear all headers and footers
\fancyhead[L]{Electrodynamics}
\fancyhead[C]{Wavefunctions}
\fancyhead[R]{Problem Set PS.18}
\fancyfoot[C]{\thepage}
\renewcommand{\headrulewidth}{0.4pt}

% Title block for the first page
\fancypagestyle{plain}{
    \fancyhf{}
    \fancyfoot[C]{\thepage}
    \renewcommand{\headrulewidth}{0pt}
}

\begin{document}
\thispagestyle{plain}

\begin{center}
    {\Large Electrodynamics}\\
    {\Large Problem Set PS.18}\\
    \vspace{0.2cm}
    Wavefunctions\\
    March 16, 2026\\
    Zoeb Izzi
\end{center}

\vspace{0.5cm}

% Problem 1
\noindent \textbf{1. Time Evolution in Infinite Square Well}
\hrule
\vspace{0.3cm}

\noindent Consider a particle of mass $m$ in a one-dimensional infinite square well extending from $x = 0$ to $x = L$ with potential:
\[ v(x) = \begin{cases} 0 & 0 < x < L \\ \infty & otherwise \end{cases} \]
We know the normalized stationary states and energy eigenvalues are given by:
\[ \phi_n(x) = \sqrt{\frac{2}{L}} \sin\left(\frac{n\pi}{L}x\right) \]
\[ E_n = \frac{n^2\pi^2\hbar^2}{2mL^2} \]
at time $t = 0$ the particle is prepared in the normalized state:
\[ \psi(x) = \begin{cases} Ax & 0 \le x \le L/2 \\ A(L - x) & L/2 \le x \le L \\ 0 & elsewhere \end{cases} \]

\begin{enumerate}[label=\alph*.]
    \item Determine the normalization constant $A$.
    \\
    \\
    \[ \int_0^L |\psi(x)|^2 dx = 1 \]
    \[ \int_0^{L/2} (Ax)^2 dx + \int_{L/2}^L (A(L-x))^2 dx = 1 \]
    By symmetry:
    \[ 2 A^2 \int_0^{L/2} x^2 dx = 1 \]
    \[ 2 A^2 \left[ \frac{x^3}{3} \right]_0^{L/2} = 2 A^2 \left( \frac{L^3/8}{3} \right) = A^2 \frac{L^3}{12} = 1 \]
    \[\therefore \hspace{1mm} \boxed{A = \sqrt{\frac{12}{L^3}} = \frac{2\sqrt{3}}{L^{3/2}}} \]

    \item We know that a general state can be written as the linear combination of eigenstates:
    \[ \psi(x) = \sum_{n=1}^\infty c_n\phi_n(x) \]
    Where
    \[ c_n = \int_0^L \phi_n^*(x)\psi(x)dx \]
    Explicitly calculate $c_1, c_2, c_3$.
    \\
    \\
    \[ c_n = \sqrt{\frac{2}{L}} A \left[ \int_0^{L/2} x \sin\left(\frac{n\pi x}{L}\right) dx + \int_{L/2}^L (L-x) \sin\left(\frac{n\pi x}{L}\right) dx \right] \]
    \[ I_1 = \left[ -\frac{Lx}{n\pi} \cos\left(\frac{n\pi x}{L}\right) + \frac{L^2}{n^2\pi^2} \sin\left(\frac{n\pi x}{L}\right) \right]_0^{L/2} \]
    \[ I_1 = -\frac{L^2}{2n\pi} \cos\left(\frac{n\pi}{2}\right) + \frac{L^2}{n^2\pi^2} \sin\left(\frac{n\pi}{2}\right) \]
    \[ \int_{L/2}^0 u \sin\left(\frac{n\pi(L-u)}{L}\right) (-du) = \int_0^{L/2} u \sin\left(n\pi - \frac{n\pi u}{L}\right) du \]
    \begin{itemize}
        \item If $n$ is \textbf{even}: $1 + (-1)^{n+1} = 0$, so $c_{even} = 0$.
        \item If $n$ is \textbf{odd}: $1 + (-1)^{n+1} = 2$. Also, for odd $n$, $\cos(n\pi/2) = 0$.
    \end{itemize}
    When $n$ is odd:
    \[ c_n = \sqrt{\frac{2}{L}} \sqrt{\frac{12}{L^3}} \times 2 \left( \frac{L^2}{n^2\pi^2} \sin\left(\frac{n\pi}{2}\right) \right) = \frac{\sqrt{24}}{L^2} \frac{2L^2}{n^2\pi^2} \sin\left(\frac{n\pi}{2}\right) = \frac{4\sqrt{6}}{n^2\pi^2} \sin\left(\frac{n\pi}{2}\right) \]
    
    \begin{align*}
        c_1 &= \frac{4\sqrt{6}}{1^2\pi^2} \sin\left(\frac{\pi}{2}\right) = \frac{4\sqrt{6}}{\pi^2} \\
        c_2 &= 0 \\
        c_3 &= \frac{4\sqrt{6}}{3^2\pi^2} \sin\left(\frac{3\pi}{2}\right) = -\frac{4\sqrt{6}}{9\pi^2}
    \end{align*}

    \item Write out the time-dependent state $\psi(x, t)$ in the form:
    \[ \psi(x, t) = \sum_{n=1}^3 c_n\phi_n(x)e^{-iE_nt/\hbar} \]
    \\
    \\
    \[ \psi(x,t) = \left(\frac{4\sqrt{6}}{\pi^2}\right) \sqrt{\frac{2}{L}} \sin\left(\frac{\pi x}{L}\right) e^{-iE_1t/\hbar} - \left(\frac{4\sqrt{6}}{9\pi^2}\right) \sqrt{\frac{2}{L}} \sin\left(\frac{3\pi x}{L}\right) e^{-iE_3t/\hbar} \]
    \[ = \frac{8\sqrt{3}}{\pi^2\sqrt{L}} \sin\left(\frac{\pi x}{L}\right) e^{-iE_1t/\hbar} - \frac{8\sqrt{3}}{9\pi^2\sqrt{L}} \sin\left(\frac{3\pi x}{L}\right) e^{-iE_3t/\hbar} \]

    \item Suppose an energy measurement is made at time $t = 0$.
    \begin{enumerate}[label=\roman*.]
        \item What are the possible results?
        \\
        \\
        The energy eigenvalues $E_n = \frac{n^2\pi^2\hbar^2}{2mL^2}$. However, since $c_n = 0$ for all even integers, it's really only $E_n$ for odd $n$.
        
        \item What is the probability of obtaining the result $E_n$?
        \\
        \[ P(E_n) = |c_n|^2 = \begin{cases} \frac{96}{n^4\pi^4} & \text{for } n \text{ odd} \\ 0 & \text{for } n \text{ even} \end{cases} \]
        
        \item What is the probability of measuring $E_1, E_2, E_3$?
        \\
        \begin{align*}
            P(E_1) &= |c_1|^2 = \frac{96}{\pi^4} \approx 0.9855 \text{ (or } 98.55\%) \\
            P(E_2) &= |c_2|^2 = 0 \\
            P(E_3) &= |c_3|^2 = \frac{96}{81\pi^4} = \frac{32}{27\pi^4} \approx 0.0122 \text{ (or } 1.22\%)
        \end{align*}
    \end{enumerate}

    \item Suppose the first energy measurement gives the value $E_3$.
    \begin{enumerate}[label=\roman*.]
        \item What is the state immediately after the measurement?
        \\
        \[ \psi_{\text{after}}(x) = \phi_3(x) = \sqrt{\frac{2}{L}} \sin\left(\frac{3\pi x}{L}\right) \]
        
        \item What is the state at a later time $t + \Delta t$?
        \\
        \[ \psi_{\text{later}}(x, t+\Delta t) = \phi_3(x) e^{-iE_3 \Delta t/\hbar} = \sqrt{\frac{2}{L}} \sin\left(\frac{3\pi x}{L}\right) e^{-iE_3 \Delta t/\hbar} \]
    \end{enumerate}

    \item Calculate $\langle H \rangle$. Should this be time-dependent?
    \\
    \\
    \[ \langle H \rangle = \sum_{n \text{ odd}} \left( \frac{96}{n^4\pi^4} \right) \left( \frac{n^2\pi^2\hbar^2}{2mL^2} \right) = \frac{48\hbar^2}{\pi^2 mL^2} \sum_{n \text{ odd}} \frac{1}{n^2} \]
    \[ \langle H \rangle = \frac{48\hbar^2}{\pi^2 mL^2} \left( \frac{\pi^2}{8} \right) = \frac{6\hbar^2}{mL^2} \]
    \textbf{No.} The potential is time-independent, which implies energy is conserved. 

    \item Calculate $\langle x \rangle$. Should this be time-dependent?
    \\
    \\
    At $t=0$, the probability distribution $|\psi(x,0)|^2$ is perfectly symmetric about the center of the well, so $\langle x \rangle = L/2$. 
        
    Therefore, $\langle x \rangle = L/2$ permanently. \textbf{No, it is not time-dependent.}

    \item What if we make a position measurement of the state $\psi(x, t)$ (not an energy measurement)? Explain the state of the system just after this measurement. Explicitly comment on its energy content and if the state is stationary or non-stationary.
    \\
    \\
    \textbf{State just after measurement:} Measuring the position collapses the wavefunction into a Dirac delta function spiked exactly at the measured position $x_0$, so $\psi_{\text{after}}(x) = \delta(x - x_0)$.
    
    \textbf{Energy content:} Heisenberg's basics tell us that knowing the position causes infinite uncertainty in momentum, which corresponds to an infinite amount of kinetic energy. $\langle H \rangle \to \infty$. 
    
    \textbf{Stationary or non-stationary?} It is \textbf{non-stationary}. The delta function is a superposition of nearly all possible energy eigenstates. As time passes, the  phase factors of different modes will rapidly interfere and reduce, leaving the wave function to take over the well.

\end{enumerate}

\end{document}